\section{The published fluid input and its complete arithmetic profile map}
\label{sec:cq-NS-profile-transfer}

We now apply the public fluid theorem recorded in
Section~\ref{sec:cq-public-NS-source} through a proved map.
The complete companion calculation \cite{cqNSFlow} was read.
The construction and every estimate used here are proved below,
including the polynomial coordinates, real flow, pressure,
forcing, physical metric, arithmetic Fourier constants, and inverse.
The final existence assertion cites the public Theorem~1.1 as
an external mathematical input; it is not an independent
verification of that theorem's complete proof.

\subsection{The original three coordinates and their full inverse differential}

In the original real coordinates $\mathbf x=(x,y,w)$ put
\begin{equation}
\begin{aligned}
 F_1&=(1+xy)^3w+y^2(1+xy)(4+3xy),\\
 F_2&=y+3x(1+xy)^2w+3xy^2(4+3xy),\\
 F_3&=2x-3x^2y-x^3w,\\
 \mathbf S&=(F_1/2,F_2,F_3)^{\mathsf T},\qquad J_S=D\mathbf S .
\end{aligned}
\label{eq:cq-NS-original-S}
\end{equation}
Thus $S_1=\sigma$ is exactly the earlier spatial coordinate.
The bold notation distinguishes this spatial three-vector
from the original quotient multiplication operator $S$.

Here is a complete determinant calculation. On $x\ne0$ set
$v=1+xy$, $h=2-3xy-x^2w$, and
$A=v^2+v-v^3h$, $B=4v+2-3v^2h$. Direct substitution gives
$F=(A/x^2,B/x,xh)$ and
\[
 \det\frac{\partial(x,v,h)}{\partial(x,y,w)}=-x^3,\qquad
 \det\frac{\partial F}{\partial(x,v,h)}
 =x^{-3}\det\begin{pmatrix}
 -2A&2v+1-3v^2h&-v^3\\
 -B&4-6vh&-3v^2\\h&0&1
 \end{pmatrix}=2x^{-3}.
\]
For the last equality the determinant's numerator is
\[
 (4-6vh)(-2v^2-2v+3v^3h)
 +(2v+1-3v^2h)(4v+2-6v^2h)=2,
\]
as expansion verifies. Thus $\det DF=-2$ on $x\ne0$ and,
by polynomial identity, everywhere. Consequently
$\det J_S=-1$ everywhere. Define the polynomial columns
\begin{equation}
\begin{aligned}
 W_1&=-\tfrac12\nabla F_2\times\nabla F_3,&
 W_2&=-\tfrac12\nabla F_3\times\nabla F_1,&
 W_3&=-\tfrac12\nabla F_1\times\nabla F_2,\\
 (DF)^{-1}&=[\,W_1\ W_2\ W_3\,],&
 J_S^{-1}&=[\,2W_1\ W_2\ W_3\,].
\end{aligned}
\label{eq:cq-NS-original-inverse}
\end{equation}
The scalar triple products with the three rows $\nabla F_i$
give $DF\,W_i=e_i$, proving both inverse identities.
In particular $\ker(d\sigma)_{\mathbf x}
=\operatorname{span}_{\mathbb R}\{W_2(\mathbf x),W_3(\mathbf x)\}$:
write any vector in the basis $(W_1,W_2,W_3)$ and use
$d\sigma(W_i)=\delta_{1i}/2$. No global inverse of $F$ or
$\mathbf S$ is presumed.

\subsection{The complete real flow and its volume-preserving inverse}

Fix $K\subset\mathbb R^3$ compact, $0<T<\infty$ and $\nu>0$.
Let $\mathcal U_{K,T,\nu}$ consist of smooth real triples
$(u,p,f)$ on $[0,T)\times\mathbb R^3$ satisfying
\begin{equation}
 \partial_tu_i+\sum_{\ell=1}^3u_\ell\partial_\ell u_i
 =\nu\Delta_{\mathbf x}u_i-\partial_i p+f_i,\qquad
 \sum_i\partial_i u_i=0,\qquad
 \operatorname{supp}u(t,\cdot)\subseteq K .
 \label{eq:cq-NS-input-class}
\end{equation}
There is no restriction on $p,f$ beyond this definition for
the correspondence; the public theorem supplies the stronger
support and initial conditions in its stated example.
For any such triple define
\[
 \partial_tX_t(q)=u(t,X_t(q)),\qquad X_0(q)=q,\qquad
 A_X(t,q)=D_qX_t(q),\qquad q\in\mathbb R^3 .
\]
For each $T_0<T$, every spatial derivative of $u$ is bounded
on $[0,T_0]\times K$ and vanishes outside $K$.
In particular its global spatial Lipschitz constant is
uniform on this time interval. On an interval of length
less than its reciprocal, the integral map
$Y\mapsto q+\int u(s,Y(s))\,ds$ is a contraction on
continuous curves in the uniform norm. Successive such
intervals construct the solution to $T_0$; boundedness of
$u$ prevents finite escape. The same construction backward
from each endpoint constructs its inverse. Differentiating
the integral equation gives a linear integral equation for
each first label derivative, and then successively for every
higher label and time derivative. The same local contraction
and induction prove smoothness and uniqueness of these
derivatives. Thus $X_t$ is a smooth global diffeomorphism
for every $t<T$.

Differentiation and the determinant product rule give
\[
 \partial_t A_X=Du(t,X_t)A_X,\quad A_X(0,q)=I_3,\qquad
 \partial_t\det A_X=(\operatorname{div}u)(t,X_t)\det A_X=0 .
\]
The determinant identity first holds where the matrix is
invertible, including a neighbourhood of $t=0$; its value
one prevents any first loss of invertibility and continues
the identity. Hence $\det A_X=1$ with positive orientation.
Every $q\notin K$ has its constant trajectory. An endpoint
outside $K$ also has a constant backward trajectory, so
inverse uniqueness proves
\begin{equation}
 X_t(K)=K,\qquad X_t(q)=q\ (q\notin K),\qquad
 X_t(L)\subseteq K\cup L\quad(L\subset\mathbb R^3\text{ compact}).
 \label{eq:cq-NS-flow-support}
\end{equation}
All these conclusions hold on each preterminal interval
without a uniform bound on $|u|$ as $t$ approaches $T$.

\subsection{The full complex fibre and both time derivatives}

Set $a_j(t,q)=S_j(X_t(q))$ and
$v_j(t,q)=(\nabla S_j\cdot u)(t,X_t(q))$.
Then $a_j$ is real and $\partial_ta_j=v_j$.
At fixed $t$ the map on the actual real and complex variables is
\[
 (q,\zeta)\longmapsto(q,a_j(t,q)+\zeta),\qquad
 (q,s)\longmapsto(q,s-a_j(t,q)).
\]
These are exact inverse maps, smooth in $q$ and holomorphic
on each complex fibre. Write
$\mathcal B=C^\infty(\mathbb R^3;\mathcal O(\mathbb C_\zeta))$
with seminorms consisting of suprema of finitely many label
and spectral derivatives on compact products.
The map
\[
 \mathcal P_t:\mathcal O(\mathbb C)^3\longrightarrow\mathcal B^3,
 \qquad \mathcal P_t(h_j)_j=(h_j(a_j(t,q)+\zeta))_j
\]
has the left inverse, for any fixed $q_*$,
$(\mathcal R_tG)_j(s)=G_j(q_*,s-a_j(t,q_*))$.
Substitution proves $\mathcal R_t\mathcal P_t=I$.
Images of compact products have compact spectral argument;
the chain rule involves only finitely many bounded derivatives
of $a_j$ there. This proves continuity of $\mathcal P_t$,
while evaluation and translation prove continuity of
$\mathcal R_t$. Therefore $\mathcal P_t$ is a topological
isomorphism onto its image with the inherited topology.
At $\zeta=0$ its first component is exactly
$h_1(\sigma(X_t(q)))=X_t^*\sigma^*h_1(q)$.

Keep a smooth, explicitly chosen real Newman-time function
$\vartheta:[0,T)\to\mathbb R$, independently of physical $t$.
For either original profile $h_\theta=Z_\theta$ or $K_\theta$,
define
\[
 \Psi_j(t,q,\zeta)=h_{\vartheta(t)}(\zeta+a_j(t,q)).
\]
The previously proved heat identity gives
\begin{equation}
 \partial_t\Psi_j
 =-\frac{\vartheta'}4\partial_\zeta^2\Psi_j
      +v_j\partial_\zeta\Psi_j .
 \label{eq:cq-NS-profile-generator}
\end{equation}
Since $v_j$ has no spectral dependence, another complete
differentiation gives
\begin{equation}
 \partial_t^2\Psi_j
 =\frac{(\vartheta')^2}{16}\partial_\zeta^4\Psi_j
 -\frac{\vartheta'v_j}{2}\partial_\zeta^3\Psi_j
 +\left(v_j^2-\frac{\vartheta''}{4}\right)\partial_\zeta^2\Psi_j
 +b_j\partial_\zeta\Psi_j,
 \label{eq:cq-NS-profile-second-generator}
\end{equation}
where, with every original momentum term,
\begin{equation}
 b_j=\partial_tv_j
 =\left[
 \sum_i(\nu\Delta u_i-\partial_i p+f_i)\partial_iS_j
 +\sum_{i,\ell}u_i u_\ell\partial_i\partial_\ell S_j
 \right](t,X_t(q)).
 \label{eq:cq-NS-profile-acceleration}
\end{equation}
Indeed the material derivative of $\nabla S_j\cdot u$
gives the Hessian term and the full material acceleration
of $u$; substituting \eqref{eq:cq-NS-input-class} gives
\eqref{eq:cq-NS-profile-acceleration}. In the second
derivative the mixed cubic term occurs twice, each with
coefficient $-\vartheta'v_j/4$, proving the factor $-1/2$.

The full Euclidean viscous operator in label coordinates is
\[
 G_X=A_X^{-1}A_X^{-\mathsf T},\qquad
 \mathcal L_t=\sum_{\alpha,\beta}
       \partial_{q_\alpha}(G_X^{\alpha\beta}\partial_{q_\beta}).
\]
For every smooth scalar $b$ one has
$\mathcal L_tX_t^*b=X_t^*\Delta b$. To prove it with all
first-order coefficient derivatives, test against
$\phi\in C_c^\infty(\mathbb R^3)$ and use the gradient chain
rule and the change of variables with determinant one:
\[
\begin{aligned}
 \int\phi\,\mathcal L_tX_t^*b\,dq
 &=-\int\nabla_q\phi\cdot A_X^{-1}(\nabla b)\circ X_t\,dq\\
 &=-\int\nabla_{\mathbf x}(\phi\circ X_t^{-1})\cdot\nabla b\,d\mathbf x
 =\int\phi\,X_t^*\Delta b\,dq .
\end{aligned}
\]
Smoothness makes this distributional identity pointwise.
The ordinary spatial chain rule now proves
\begin{equation}
\begin{aligned}
 \mathcal L_t\Psi_j
 &=(|\nabla S_j|^2\circ X_t)\partial_\zeta^2\Psi_j
        +(\Delta S_j\circ X_t)\partial_\zeta\Psi_j,\\
 (\partial_t-\nu\mathcal L_t)\Psi_j
 &=-\left(\frac{\vartheta'}4+\nu|\nabla S_j|^2\circ X_t\right)
          \partial_\zeta^2\Psi_j
       +(v_j-\nu\Delta S_j\circ X_t)\partial_\zeta\Psi_j .
\end{aligned}
\label{eq:cq-NS-material-viscous-residual}
\end{equation}

\subsection{Original Fourier constants and exact velocity recovery}

Define, with the original arithmetic kernel,
\[
 h_*(r)=\frac{\pi}{2}r^2(2\pi r^2-3)e^{-\pi r^2},\qquad
 k(\omega)=e^{\omega/2}\sum_{n\geq1}h_*(ne^\omega).
\]
For the Fourier convention with exponential $e^{-2\pi ir\xi}$,
the second and fourth Gaussian derivative formulas are
\[
 \widehat{r^2e^{-\pi r^2}}
   =(1/(2\pi)-\xi^2)e^{-\pi\xi^2},\quad
 \widehat{r^4e^{-\pi r^2}}
   =(\xi^4-3\xi^2/\pi+3/(4\pi^2))e^{-\pi\xi^2}.
\]
They follow by differentiating the Gaussian transform twice
and four times and dividing by $(-2\pi i)^2$ and
$(-2\pi i)^4$, respectively. The coefficients
$\pi^2$ and $-3\pi/2$ then give $\widehat h_*=h_*$ and
$h_*(0)=0$. Poisson summation gives
$\sum_{n\geq1}h_*(n\lambda)
=\lambda^{-1}\sum_{n\geq1}h_*(n/\lambda)$, and hence
$k(\omega)=k(-\omega)$. Poisson summation for this Schwartz
function follows directly by periodizing it: its Fourier
coefficients are the displayed Fourier transform at integer
frequencies; the function and coefficient series converge
absolutely with derivatives, so evaluation at zero is valid.

Each fixed number of derivatives of the defining series is
locally uniformly convergent. At $\omega\geq0$ its Gaussian
terms give
$|k^{(r)}(\omega)|\leq C_r e^{C_r\omega}
 e^{-\pi e^{2\omega}}$; factor the $n=1$ exponential and
bound the remaining polynomially weighted series by
$\sum n^{M_r}e^{-\pi(n^2-1)}<\infty$.
Evenness gives the same bound with $|\omega|$ at the other end.
Expanding the definition proves
$\Phi(u)=2k(2u)$ with every factor retained. Consequently
\begin{equation}
\begin{aligned}
 Z_\theta(s)&=4\int_{\mathbb R}k(\omega)
                       e^{\theta\omega^2/4}e^{is\omega}\,d\omega,\\
 K_\theta(s)&=4\int_{\mathbb R}m(\omega)k(\omega)
                       e^{\theta\omega^2/4}e^{is\omega}\,d\omega ,
\end{aligned}
\label{eq:cq-NS-original-Fourier}
\end{equation}
where $m$ is exactly the three-shift, four-term multiplier
in \eqref{eq:cq-xi4-full-Fourier-kernel}, including
$-i\omega/336$ at shift $3$.
All derivatives of these amplitudes are Schwartz for real
$\theta$, locally uniformly in $\theta$; the same is true
of their Fourier transforms by integration by parts.

Set $d_h(\theta)^2=\int_{\mathbb R}|h_\theta'(s)|^2\,ds$.
The exact constants are
\begin{equation}
\begin{aligned}
 d_Z(\theta)^2&=32\pi\int_{\mathbb R}\omega^2 k(\omega)^2
                                   e^{\theta\omega^2/2}\,d\omega,\\
 d_K(\theta)^2&=32\pi\int_{\mathbb R}\omega^2|m(\omega)|^2
                            k(\omega)^2e^{\theta\omega^2/2}\,d\omega .
\end{aligned}
\label{eq:cq-NS-profile-norm-constants}
\end{equation}
For completeness, insert $e^{-\epsilon s^2}$ in the
squared norm of $4\int b(\omega)e^{is\omega}\,d\omega$.
Fubini and the Gaussian integral give
\[
 16\sqrt{\pi/\epsilon}\iint
 b(\omega)\overline{b(\eta)}
             e^{-(\omega-\eta)^2/(4\epsilon)}\,d\omega\,d\eta .
\]
The Gaussian convolution has mass $2\pi$ and tends to
$2\pi b$ in this integral: split at fixed distance from
zero, use uniform continuity on the near part and its
Gaussian tail on the far part, with Schwartz domination.
On the original side dominated convergence applies because
the Fourier transform is Schwartz. The result is
$32\pi\int|b|^2$. Taking
$b=i\omega k e^{\theta\omega^2/4}$ or
$i\omega m k e^{\theta\omega^2/4}$ proves
\eqref{eq:cq-NS-profile-norm-constants}.
Both integrals are strictly positive:
$k(0)=\Phi(0)/2>0$ and $m(0)=127/219024>0$
give an interval containing nonzero frequencies where the
integrands' non-Gaussian factors are positive.
The two-ended bounds prove continuity in $\theta$ by
domination on compact parameter intervals. Thus on each
compact $I\subset\mathbb R$ there are
$0<d_-\leq d_h(\theta)\leq d_+<\infty$.

Define the full arithmetic temporal defect on the real
spectral line by
\[
 \mathscr D_j=\partial_t\Psi_j+
                \frac{\vartheta'}4\partial_s^2\Psi_j
        =v_j h_{\vartheta}'(s+a_j).
\]
Translation by the real $a_j$ preserves Lebesgue measure.
Therefore
\begin{equation}
\begin{aligned}
 \sum_j\|\mathscr D_j(t,q,\cdot)\|_2^2
       &=d_h(\vartheta(t))^2|v(t,q)|^2,\\
 v_j(t,q)&=\frac{\int_{\mathbb R}
       \overline{h_{\vartheta}'(s+a_j)}\,\mathscr D_j(t,q,s)\,ds}
                {d_h(\vartheta)^2},\\
 u(t,X_t(q))&=2W_1(X_t(q))v_1+
                    W_2(X_t(q))v_2+W_3(X_t(q))v_3 .
\end{aligned}
\label{eq:cq-NS-exact-velocity-recovery}
\end{equation}
These formulas prove an inverse correspondence
$(X_t,u\circ X_t)\leftrightarrow(X_t,\mathscr D)$ on the
three real profile lines at each $(t,q)$. The flow and its
inverse sheet are retained in both sides.

For nonempty $K$, put
$M_J=\max_K\|J_S\|_{\rm op}$ and
$N_J=\max_K\|J_S^{-1}\|_{\rm op}$.
They are finite and positive. Define
\[
 \mathcal A_\infty(t)=
 \sup_q\left(\sum_j\|\mathscr D_j(t,q,\cdot)\|_2^2\right)^{1/2},
 \quad
 \mathcal A_2(t)=
 \left(\sum_j\int_{\mathbb R^3}\int_{\mathbb R}
               |\mathscr D_j|^2\,ds\,dq\right)^{1/2}.
\]
When $\vartheta([0,T))\subseteq I$, the complete estimates are
\begin{equation}
 \frac{d_-}{N_J}\|u(t)\|_\infty\leq\mathcal A_\infty(t)
       \leq d_+M_J\|u(t)\|_\infty,\qquad
 \frac{d_-}{N_J}\|u(t)\|_2\leq\mathcal A_2(t)
       \leq d_+M_J\|u(t)\|_2 .
 \label{eq:cq-NS-energy-and-supremum}
\end{equation}
Indeed $v=J_S(X_t)(u\circ X_t)$, both vanish off $K$,
and $|J_Su|\leq M_J|u|$, $|u|\leq N_J|J_Su|$ on $K$.
The supremum is preserved by the diffeomorphism $X_t$.
For the integral, its determinant one gives the exact identity
\[
 \mathcal A_2(t)^2
 =d_h(\vartheta(t))^2
           \int_{\mathbb R^3}|J_S(\mathbf x)u(t,\mathbf x)|^2\,d\mathbf x .
\]
This proves both bounds without discarding any spatial or
spectral component. For a fixed original incidence coefficient
$a_{ef}$, the observation $a_{ef}K_\theta$ multiplies
$\mathscr D^K$ by $a_{ef}$ and its squared norm by
$|a_{ef}|^2$; inversion uses its actual nonzero value, and
the zero incidence has exactly the zero observation.

\subsection{Application to the stated public singularity theorem}

\begin{cqtheorem}[Arithmetic consequence of the cited public input]
\label{thm:cq-NS-public-propagation}
Take the velocity, pressure, force and compact set supplied
by Theorem~1.1 of \cite{cqNS2026}, with its original $\nu>0$.
Choose the explicit Newman time $\vartheta(t)=0$ and the
actual arithmetic profile $h_0=\Xi$, on $0\leq t<1$.
Then the constructed entire spectral fields satisfy
\begin{equation}
 \mathscr D_j=\partial_t\Psi_j,\qquad
 \mathscr D_j(0,q,s)=0,\qquad
 \sup_{t<1}\mathcal A_2(t)<\infty,\qquad
 \limsup_{t\uparrow1}\mathcal A_\infty(t)=\infty .
 \label{eq:cq-NS-public-arithmetic-blowup}
\end{equation}
Every $\mathscr D_j$ is supported in $q\in K$.
For every compact $L\subset\mathbb R^3$, compact
$E_0\subset\mathbb C$ and integer $r\geq0$,
\begin{equation}
 \sup_{t<1,\ q\in L,\ \zeta\in E_0}
             |\partial_\zeta^r\Psi_j(t,q,\zeta)|<\infty .
 \label{eq:cq-NS-bounded-spectral-profiles}
\end{equation}
The analogous temporal-defect conclusions hold for the
original full profile $K_0$ with its exact $d_K(0)$.
\end{cqtheorem}

\begin{proof}
The cited public theorem provides a triple in
\eqref{eq:cq-NS-input-class} with $T=1$, $u(0)=0$,
uniformly bounded $\|u(t)\|_2$, and unbounded terminal
$\|u(t)\|_\infty$. The preceding explicit construction
therefore supplies $X_t,\Psi,\mathscr D$. With the chosen
constant Newman time the heat part of the temporal defect
vanishes exactly, and $u(0)=0$ gives $\mathscr D(0)=0$.
Set $d_-=d_+=d_Z(0)>0$ in
\eqref{eq:cq-NS-energy-and-supremum}. Its two inequalities
give the energy bound and the unbounded terminal supremum
in \eqref{eq:cq-NS-public-arithmetic-blowup}.
Support follows from $v=0$ off $K$.
For the last bound,
\eqref{eq:cq-NS-flow-support} places the argument of
$\Xi^{(r)}$ in the fixed compact set $S_j(K\cup L)+E_0$.
Entire continuity gives its finite maximum.
All arguments apply to $K_0$ using the already proved
positive constant $d_K(0)$ and its entire derivatives.
This proof uses the external existence theorem exactly at
its first step; all transfer formulas and estimates are
proved within this fragment.
\end{proof}
